\scanpage{20}
\mspage{19}

- 15 -\hfill n° 391

d'où $s'(x)=s(x)+2t$, et en remplaçant $x$ par $s(x)$, il vient $s'(s(x))=x+2t$.

Dans le cas particulier où $K=R$ et où $\unsure{\underline{\Phi}|\Phi}$ est positive non dégénérée, on vérifie aussitôt que $\|t\|$ est la \uline{distance} des hyperplans $H$, $H'$ (pour la structure d'espace métrique définie sur $E$ par la distance euclidienne).

5. \uline{Chambres dans un espace affine.}

Soient $E$ un \uline{espace affine réel} de dimension finie $n$ ; en prenant une origine dans $E$, $E$ s'identifie à un espace vectoriel réel de dimension $n$ et est donc muni d'une topologie bien déterminée pour lequel il est isomorphe à $R^n$ (\uline{Top.} \uline{gén.}, chap. VI, §1, n° 5) ; en outre, les translations étant des homéomorphismes de $R^n$ sur lui-même, cette topologie est indépendante de l'origine choisie dans $E$ ; quand nous considérerons dans ce qui suit $E$ comme espace topologique, il s'agira toujours de cette topologie.

Pour tout hyperplan $H$ dans $E$ et tout ensemble \uline{connexe} non vide $C$ ne rencontrant pas $H$, $C$ est contenu dans un des demi-espaces ouverts déterminés par $H$, ces derniers formant une partition de $E-H$ en deux ensembles ouverts connexes ; on désignera le demi-espace ouvert déterminé par $H$ et contenant $C$ par $D_H(C)$.

LEMME 1.- \uline{Si $\underline H$ est un ensemble localement fini} (\uline{Top.} \uline{gén.}, chap. I, 3\textsuperscript{e} éd., §1, n° 5, déf. 6) \uline{d'hyperplans affines de $E$, la réunion} $S'$ \uline{des} $H\in\underline H$ \uline{est un ensemble fermé et rare dans $E$, et pour toute composante connexe} $C$ \uline{de} $E-S$, \uline{on a} $C=\bigcap_{H\in\underline H}D_H(C)$.

Comme les $H\in\underline H$ sont fermés, on sait qu'il en est de même de $S$ (\uline{Top.} \uline{gén.}, chap. I, 3\textsuperscript{e} éd., §1, n° 5, prop. 4) ; d'autre part, chacun des $H\in\underline H$ est rare dans $E$ ; si $x\in E$, il y a un voisinage ouvert $V$ de $x$ ne rencontrant qu'un nombre fini d'éléments $H_i\in\underline H$ ($1\leq i\leq n$), et la réunion des $V\cap H_i$ est rare dans $V$, puisqu'il en est ainsi de chacun d'eux. Pour tout $H\in\underline H$, il est clair
